% \setcounter{page}{1}
\section{Introduction}

The number of electric vehicles and charging infrastructure in Germany has grown substantially in recent years. Today, there are more than 1.6 million battery-operated vehicles on German roads, with more joining every day. Alongside the increase in EVs, the charging network is becoming denser too, ranging from fast charging stations on motorways to \ac{AC} stations in city centers and shopping centers \cite{noauthor_elektroauto-statistiken_nodate}. 

However, simply increasing the number of charging stations will not be enough. The locations and integration into the power grid must be thoroughly planned to, for example, balance peak and off-peak loads and optimize user needs.

The German government plans to have around 15 million electric vehicles on German roads by 2030, and aims to provide around one million public charging stations by then \cite{klimaschutz_rahmenbedingungen_nodate}. To actually realize that economically and ecologically, concepts like intelligent locations are becoming more important.

The next step in innovation would be bidirectional charging (\ac{V2G}). This would mean that electric vehicles would not only be a means of transportation, but also a form of mobile energy storage that could feed surplus power back into the grid. This would help to cushion load peaks and increase the stability of the overall system, making an important contribution to the energy transition, especially in combination with volatile renewables.

These developments clearly show that simply increasing the number of charging stations is not enough. It is becoming increasingly important to integrate these technologies in a smart way in order to promote electric mobility sustainably and across the board.

\subsection{Motivation}
Reliably testing and validating the positions of charging stations in a controlled environment would be a significant step toward an automated ecosystem in which companies or government entities could respond to demand. In short, the goal is to develop a system that improves overall efficiency and effectiveness while producing meaningful and realistic results.

\subsection{Relevance}

Simulation tools can be used to recreate realistic traffic and charging behavior in a virtual environment, enabling testing different locations and operational strategies at an early stage. By modeling driving patterns and charging workflows, it is possible to analyze typical routes and charging patterns. This allows charging stations to be placed where they will be most beneficial and minimize no-load times. In addition, grid loads can be estimated in advance to avoid overloads and enable expansion where the power grid still has capacity.

In terms of bidirectional charging, simulation tools can help combine driver behavior and network analysis. This makes it possible to calculate how much energy can be fed back into the grid without creating voltage peaks. Different scenarios can be simulated, such as the bulk loading that occurs when people get home from work and connect their cars to their home charging stations.

Using simulation to support planning processes can ensure greater reliability and efficiency by providing data-based insights into optimal locations, performance, and network requirements. Eliminating the need for on-site tests and pilot systems can significantly reduce planning and testing costs. Simulation forms the basis of smart, well-thought-out charging infrastructure.